\documentclass[11pt,a4paper]{article}
\usepackage[a4paper,top=3.6cm,bottom=2.4cm,left=1.8cm,right=1.8cm,headheight=44pt,headsep=14pt]{geometry}
\usepackage{amsmath,amssymb,mathtools}
\usepackage{graphicx}
\usepackage{xcolor}
\usepackage{array,tabularx,longtable,booktabs}
\usepackage{fancyhdr}
\usepackage{lastpage}
\usepackage{needspace}
\usepackage{multicol}
\setlength{\columnsep}{0.9cm}\setlength{\columnseprule}{0.4pt}
\setlength{\parindent}{0pt}
\setlength{\emergencystretch}{3em}
\definecolor{schoolblue}{RGB}{20,50,120}
\pagestyle{fancy}
\fancyhf{}
\renewcommand{\headrulewidth}{0.8pt}
\renewcommand{\footrulewidth}{0.5pt}
\fancyhead[C]{\textcolor{schoolblue}{\Large\bfseries SAGAR PUBLIC SCHOOL}\\[2pt]
  \small\textbf{Integrated Programme}\ \ $\mid$\ \ \textbf{Topic: Electrostatics -- Conductors, Gauss's Law \& Plates}}
\fancyfoot[L]{\small Sagar Public School -- Integrated Programme}
\fancyfoot[C]{\small Page \thepage\ of \pageref{LastPage}}
\fancyfoot[R]{\small Practice Sheet 3\ $\bullet$\ Electrostatics}
\fancypagestyle{plain}{\fancyhf{}\renewcommand{\headrulewidth}{0.8pt}\renewcommand{\footrulewidth}{0.5pt}%
  \fancyhead[C]{\textcolor{schoolblue}{\Large\bfseries SAGAR PUBLIC SCHOOL}\\[2pt]\small\textbf{Integrated Programme}\ \ $\mid$\ \ \textbf{Topic: Electrostatics -- Conductors, Gauss's Law \& Plates}}%
  \fancyfoot[L]{\small Sagar Public School -- Integrated Programme}\fancyfoot[C]{\small Page \thepage\ of \pageref{LastPage}}\fancyfoot[R]{\small Practice Sheet 3\ $\bullet$\ Electrostatics}}
\newcounter{qn}
\newcommand{\Q}{\needspace{5\baselineskip}\refstepcounter{qn}\par\medskip\noindent\textbf{Q.\theqn}\hspace{0.5em}}
\newcommand{\Qf}{\needspace{7.5cm}\refstepcounter{qn}\par\medskip\noindent\textbf{Q.\theqn}\hspace{0.5em}}
\newcommand{\lab}[1]{(\ifcase#1\or A\or B\or C\or D\fi)}
\newcommand{\optsfour}[4]{\par\smallskip\noindent\begin{tabularx}{\linewidth}{@{}*{4}{>{\raggedright\arraybackslash}X}@{}}
  \lab1\ #1 & \lab2\ #2 & \lab3\ #3 & \lab4\ #4\end{tabularx}\par}
\newcommand{\optstwo}[4]{\par\smallskip\noindent\begin{tabularx}{\linewidth}{@{}*{2}{>{\raggedright\arraybackslash}X}@{}}
  \lab1\ #1 & \lab2\ #2 \\[4pt] \lab3\ #3 & \lab4\ #4\end{tabularx}\par}
\newcommand{\optsv}[4]{\par\smallskip\begin{list}{}{\setlength{\leftmargin}{1.7em}\setlength{\itemsep}{1pt}\setlength{\topsep}{0pt}\setlength{\parsep}{0pt}}
  \item[\lab1] #1 \item[\lab2] #2 \item[\lab3] #3 \item[\lab4] #4\end{list}}
\newcommand{\fig}[1]{\par\smallskip\begin{center}\includegraphics[width=\linewidth,height=4.6cm,keepaspectratio]{figs/#1.png}\end{center}\par}
\newcommand{\etag}[1]{\ {\footnotesize\textit{[#1]}}}
\begin{document}
\begin{center}
{\Large\bfseries PRACTICE SHEET 3}\\[3pt]
{\normalsize\bfseries Conductors \& Charge Distribution $\bullet$ Gauss's Law \& Electric Flux $\bullet$ Electric Field \& Potential due to Sheets/Plates}
\end{center}
\vspace{2pt}
\noindent\fbox{\parbox{\dimexpr\linewidth-2\fboxsep-2\fboxrule\relax}{%
\textbf{Time Allowed: 60 Mins}\hfill \textbf{Total Questions: 20}\\[5pt]
Name: \underline{\hspace{6.3cm}}\hfill Class/Sec: \underline{\hspace{2.2cm}}\hfill Date: \underline{\hspace{2.6cm}}\\[5pt]
{\small\textbf{Instructions:} Attempt all questions. Each question carries one correct option unless marked \emph{one or more than one option is correct}. Questions without options are numerical-value type: write the answer in the blank. The tag after each question shows its source examination. A self-assessment page follows; the answer key is provided as a separate sheet.}}}
\bigskip
\begin{multicols}{2}
\Q Two metal plates (A, B) are kept horizontally with separation of $\left(\dfrac{12}{\pi}\right)$ cm, with plate A on the top. An atomizer jet sprays oil (density $1.5\ \mathrm{g/cm^3}$) droplets of radius $1$ mm horizontally. All oil droplets carry a charge $5$ nC. The potentials $V_A$ and $V_B$ are required on plates A and B respectively in order to ensure the droplets do not descend. The values of $V_A$ and $V_B$ are \underline{\hspace{1.2cm}}. (Neglect the air resistance to the droplets and take $g=10\ \mathrm{m/s^2}$)\etag{JEE Main 2026, 2 Apr, Eve.}
\optstwo{$100$ V and $580$ V}{$580$ V and $100$ V}{$60$ V and $400$ V}{$0$ V and $-200$ V}

\Qf Consider two infinitely large plane parallel conducting plates as shown below. The plates are uniformly charged with a surface charge density $+\sigma$ and $-2\sigma$. The force experienced by a point charge $+q$ placed at the mid point between two plates will be:\etag{JEE Main 2025, 2 Apr, Mor.}
\fig{f39}
\optstwo{$\dfrac{3\sigma q}{4\epsilon_0}$}{$\dfrac{3\sigma q}{2\epsilon_0}$}{$\dfrac{\sigma q}{4\epsilon_0}$}{$\dfrac{\sigma q}{2\epsilon_0}$}

\Qf A uniform electric field $E=\left(8\,\dfrac{m}{e}\right)$ V/m is created between two parallel plates of length $1$ m as shown in figure, (where $m$ = mass of electron and $e$ = charge of electron). An electron enters the field symmetrically between the plates with a speed of $2$ m/s. The angle of the deviation $(\theta)$ of the path of the electron as it comes out of the field will be \underline{\hspace{1.2cm}}.\etag{JEE Main 2022, 28 Jul, Eve.}
\fig{f96}
\optstwo{$\tan^{-1}(4)$}{$\tan^{-1}(2)$}{$\tan^{-1}\!\left(\frac13\right)$}{$\tan^{-1}(3)$}

\Q A uniform electric field $\vec E$ exists between the plates of a charged condenser. A charged particle enters the space between the plates and perpendicular to $\vec E$. The path of the particle between the plates is a:\etag{JEE Main 2013, 9 Apr, Mor.}
\optstwo{straight line}{hyperbola}{parabola}{circle}

\Qf A uniform electric field of $10$ N/C is created between two parallel charged plates (as shown in figure). An electron enters the field symmetrically between the plates with a kinetic energy $0.5$ eV. The length of each plate is $10$ cm. The angle $(\theta)$ of deviation of the path of electron as it comes out of the field is \underline{\hspace{1.5cm}} (in degree).\etag{JEE Main 2023, 25 Jan, Mor.}
\fig{f280}

\Q Two metal spheres of radius $R$ and $3R$ have same surface charge density $\sigma$. If they are brought in contact and then separated, the surface charge density on smaller and bigger sphere becomes $\sigma_1$ and $\sigma_2$, respectively. The ratio $\dfrac{\sigma_1}{\sigma_2}$ is\etag{JEE Main 2025, 8 Apr, Eve.}
\optsfour{$\dfrac13$}{$\dfrac19$}{$9$}{$3$}

\Q A small uncharged conducting sphere is placed in contact with an identical sphere but having $4\times10^{-8}$ C charge and then removed to a distance such that the force of repulsion between them is $9\times10^{-3}$ N. The distance between them is (Take $\dfrac{1}{4\pi\epsilon_0}$ as $9\times10^{9}$ in SI units)\etag{JEE Main 2025, 24 Jan, Eve.}
\optsfour{$1$ cm}{$2$ cm}{$4$ cm}{$3$ cm}

\Q Two identical conducting spheres P and S with charge $Q$ on each, repel each other with a force $16$ N. A third identical uncharged conducting sphere R is successively brought in contact with the two spheres. The new force of repulsion between P and S is:\etag{JEE Main 2024, 6 Apr, Eve.}
\optsfour{$1$ N}{$6$ N}{$12$ N}{$4$ N}

\Q The vehicles carrying inflammable fluids usually have metallic chains touching the ground:\etag{JEE Main 2024, 5 Apr, Eve.}
\optsv{To protect tyres from catching dirt from ground}{It is a custom}{To alert other vehicles}{To conduct excess charge due to air friction to ground and prevent sparking}

\Q Two isolated metallic solid spheres of radii $R$ and $2R$ are charged such that both have same charge density $\sigma$. The spheres are then connected by a thin conducting wire. If the new charge density of the bigger sphere is $\sigma'$. The ratio $\dfrac{\sigma'}{\sigma}$ is:\etag{JEE Main 2023, 30 Jan, Mor.}
\optsfour{$\dfrac53$}{$\dfrac56$}{$\dfrac94$}{$\dfrac43$}

\Q Two isolated conducting spheres $S_1$ and $S_2$ of radius $\frac23R$ and $\frac13R$ have $12\,\mu$C and $-3\,\mu$C charges, respectively, and are at a large distance from each other. They are now connected by a conducting wire. A long time after this is done the charges on $S_1$ and $S_2$ are respectively:\etag{JEE Main 2020, 3 Sep, Mor.}
\optstwo{$4.5\,\mu$C on both}{$+4.5\,\mu$C and $-4.5\,\mu$C}{$6\,\mu$C and $3\,\mu$C}{$3\,\mu$C and $6\,\mu$C}

\Q Consider a finite insulated, uncharged conductor placed near a finite positively charged conductor. The uncharged body must have a potential:\etag{JEE Main 2013, 23 Apr, Mor.}
\optsv{less than the charged conductor and more than at infinity.}{more than the charged conductor and less than at infinity.}{more than the charged conductor and more than at infinity.}{less than the charged conductor and less than at infinity.}

\Q A point charge causes an electric flux of $-2\times10^{4}\ \mathrm{Nm^2C^{-1}}$ to pass through a spherical Gaussian surface of $8.0$ cm radius, centred on the charge. The value of the point charge is: (Given $\epsilon_0=8.85\times10^{-12}\ \mathrm{C^2N^{-1}m^{-2}}$)\etag{JEE Main 2025, 29 Jan, Eve.}
\optstwo{$17.7\times10^{-8}$ C}{$-17.7\times10^{-8}$ C}{$15.7\times10^{-8}$ C}{$-15.7\times10^{-8}$ C}

\Qf $C_1$ and $C_2$ are two hollow concentric cubes enclosing charges $2Q$ and $3Q$ respectively as shown in figure. The ratio of electric flux passing through $C_1$ and $C_2$ is:\etag{JEE Main 2024, 1 Feb, Eve.}
\fig{f64}
\optsfour{$3:2$}{$5:2$}{$2:5$}{$2:3$}

\Q An electric field is given by $(6\hat i+5\hat j+3\hat k)$ N/C. The electric flux through a surface area $30\,\hat i\ \mathrm{m^2}$ lying in YZ-plane (in SI unit) is:\etag{JEE Main 2024, 29 Jan, Eve.}
\optsfour{$60$}{$90$}{$180$}{$150$}

\Q In a cuboid of dimension $2L\times2L\times L$, a charge $q$ is placed at the center of the surface $S$ having area of $4L^2$. The flux through the opposite surface to $S$ is given by\etag{JEE Main 2023, 29 Jan, Mor.}
\optstwo{$\dfrac{q}{2\epsilon_0}$}{$\dfrac{q}{3\epsilon_0}$}{$\dfrac{q}{12\epsilon_0}$}{$\dfrac{q}{6\epsilon_0}$}

\Q Two charges of $5Q$ and $-2Q$ are situated at the points $(3a,0)$ and $(-5a,0)$ respectively. The electric flux through a sphere of radius $4a$ having center at origin is:\etag{JEE Main 2024, 29 Jan, Mor.}
\optstwo{$\dfrac{2Q}{\epsilon_0}$}{$\dfrac{7Q}{\epsilon_0}$}{$\dfrac{3Q}{\epsilon_0}$}{$\dfrac{5Q}{\epsilon_0}$}

\Q Choose the incorrect statement:\par (1) The electric lines of force entering into a Gaussian surface provide negative flux.\par (2) A charge $q$ is placed at the centre of a cube. The flux through all the faces will be the same.\par (3) In a uniform electric field net flux through a closed Gaussian surface containing no net charge, is zero.\par (4) When electric field is parallel to a Gaussian surface, it provides a finite non-zero flux.\par Choose the most appropriate answer from the options given below:\etag{JEE Main 2021, 31 Aug, Eve.}
\optstwo{(3) and (4) only}{(2) and (4) only}{(4) only}{(1) and (3) only}

\Q In finding the electric field using Gauss Law the formula $|\vec E|=\dfrac{q_{enc}}{\epsilon_0|A|}$ is applicable. In the formula $\epsilon_0$ is permittivity of free space, $A$ is the area of Gaussian surface and $q_{enc}$ is charge enclosed by the Gaussian surface. The equation can be used in which of the following situation?\etag{JEE Main 2020, 8 Jan, Mor.}
\optsv{Only when $|\vec E|=$ constant on the surface.}{For any choice of Gaussian surface.}{Only when the Gaussian surface is an equipotential surface.}{Only when the Gaussian surface is an equipotential surface and $|\vec E|$ is constant on the surface.}

\Q If the electric flux entering and leaving an enclosed surface respectively is $\phi_1$ and $\phi_2$, the electric charge inside the surface will be\etag{AIEEE 2003}
\optstwo{$(\phi_2-\phi_1)\epsilon_0$}{$(\phi_2+\phi_1)/\epsilon_0$}{$(\phi_1-\phi_2)/\epsilon_0$}{$(\phi_1+\phi_2)\epsilon_0$}
\end{multicols}

\clearpage
\begin{center}{\Large\bfseries SELF-ASSESSMENT -- PRACTICE SHEET 3}\end{center}
\medskip
\renewcommand{\arraystretch}{1.9}
\noindent\begin{tabularx}{\linewidth}{|X|X|}\hline
\textbf{Time Allowed:} 60 Mins & \textbf{Time Taken:} \underline{\hspace{2.4cm}} Mins\\ \hline
\textbf{Total Questions:} 20 & \textbf{Questions Attempted:} \underline{\hspace{2cm}}\\ \hline
\textbf{Correct:} \underline{\hspace{2.6cm}} & \textbf{Incorrect:} \underline{\hspace{2.6cm}}\\ \hline
\textbf{Unattempted:} \underline{\hspace{2.2cm}} & \textbf{Accuracy:} \underline{\hspace{1.6cm}}\,\% \ {\footnotesize(Correct$\,/\,$Attempted$\,\times100$)}\\ \hline
\end{tabularx}
\renewcommand{\arraystretch}{1.0}
\bigskip
\textbf{Question-wise record} \ {\small(tick one: C = Correct, W = Wrong, S = Skipped)}\par\smallskip
{\small\renewcommand{\arraystretch}{1.55}
\noindent\begin{tabularx}{\linewidth}{|*{4}{>{\centering\arraybackslash}X|}}\hline
\textbf{Q.1}\ $\square$C\ $\square$W\ $\square$S & \textbf{Q.2}\ $\square$C\ $\square$W\ $\square$S & \textbf{Q.3}\ $\square$C\ $\square$W\ $\square$S & \textbf{Q.4}\ $\square$C\ $\square$W\ $\square$S \\ \hline
\textbf{Q.5}\ $\square$C\ $\square$W\ $\square$S & \textbf{Q.6}\ $\square$C\ $\square$W\ $\square$S & \textbf{Q.7}\ $\square$C\ $\square$W\ $\square$S & \textbf{Q.8}\ $\square$C\ $\square$W\ $\square$S \\ \hline
\textbf{Q.9}\ $\square$C\ $\square$W\ $\square$S & \textbf{Q.10}\ $\square$C\ $\square$W\ $\square$S & \textbf{Q.11}\ $\square$C\ $\square$W\ $\square$S & \textbf{Q.12}\ $\square$C\ $\square$W\ $\square$S \\ \hline
\textbf{Q.13}\ $\square$C\ $\square$W\ $\square$S & \textbf{Q.14}\ $\square$C\ $\square$W\ $\square$S & \textbf{Q.15}\ $\square$C\ $\square$W\ $\square$S & \textbf{Q.16}\ $\square$C\ $\square$W\ $\square$S \\ \hline
\textbf{Q.17}\ $\square$C\ $\square$W\ $\square$S & \textbf{Q.18}\ $\square$C\ $\square$W\ $\square$S & \textbf{Q.19}\ $\square$C\ $\square$W\ $\square$S & \textbf{Q.20}\ $\square$C\ $\square$W\ $\square$S \\ \hline
\end{tabularx}}
\bigskip
\textbf{Reason for mistakes} \ {\small(tick all that apply)}\par\smallskip
\begin{tabularx}{\linewidth}{@{}XXX@{}}$\square$ Concept not clear & $\square$ Formula forgotten & $\square$ Calculation mistake\\[3pt] $\square$ Misread the question & $\square$ Ran out of time & $\square$ Guess work\end{tabularx}
\bigskip
\textbf{Personal Feedback}\par\smallskip
\textit{What went well:}\par\vspace{2pt}
\noindent\rule{0pt}{1.35em}\hrulefill\\ \rule{0pt}{1.35em}\hrulefill\\[4pt]
\textit{Topics / questions I need to revise:}\par\vspace{2pt}
\noindent\rule{0pt}{1.35em}\hrulefill\\ \rule{0pt}{1.35em}\hrulefill\\[4pt]
\textit{My action plan before the next sheet:}\par\vspace{2pt}
\noindent\rule{0pt}{1.35em}\hrulefill\\ \rule{0pt}{1.35em}\hrulefill\\[8pt]
\textit{Teacher's remarks:}\par\vspace{2pt}
\noindent\rule{0pt}{1.35em}\hrulefill\\ \rule{0pt}{1.35em}\hrulefill\\[14pt]
Student's signature: \underline{\hspace{4cm}}\hfill Teacher's signature: \underline{\hspace{4cm}}
\end{document}
